\section{The Hecke trace limit and the principal polar coordinates}
\label{sec:hecke-assembly}
The local limit and the uniform frequency bound now give the trace
formula for the deep Hankel selection. This section completes its
principal sector by proving that the two missing coordinates become
vertical in the enlarged Hilbert space.

\begin{proposition}[The principal Hecke graph]
\label{prop:hecke-principal-graph}
Let $\Pi_\Lambda=\Pi_{\Lambda,56}$ for the principal character of
$K=\Q(\sqrt{-3})$. Let $Q_\Lambda$ project onto
$E_1(\operatorname{Ran}\Pi_\Lambda\cap\ker I\cap\ker\operatorname{ev}_0)$.
On $L^2(\R)\oplus\C^2$ define
\[
 J\psi=(E_1\psi,\psi(0),I\psi),\qquad
 \mathcal Q_\Lambda=P_{J\operatorname{Ran}\Pi_\Lambda}.
\]
For all sufficiently large $\Lambda$, the constraint map has rank two,
and
\begin{equation}\label{eq:hecke-vertical-graph}
 \bigl\|\mathcal Q_\Lambda-(Q_\Lambda\oplus0)
                         -P_{\C^2}\bigr\|_1
                         =O(\Lambda^{-1/4}).
\end{equation}
For $H(z)=\int h(t)e^{zt}dt$ and
$R(h)=V(h)\oplus H(-1/2)\oplus H(1/2)$, it follows that
\[
 \Tr(\mathcal Q_\Lambda R(h))-
           \Tr(Q_\Lambda V(h))
                  \longrightarrow H(-1/2)+H(1/2).
\]
No hypothesis on the zeros is used.
\end{proposition}
\begin{proof}
Choose $q,r\in C_c^\infty([0,\infty))$ with
\[
 q(0)=1,\quad Iq=0,\qquad r(0)=0,\quad Ir=1.
\]
Take $r$ to be a normalized bump away from zero, and obtain $q$
by subtracting a multiple of $r$ from a bump equal to one near zero.
The functions $q(|z|^2)$ and $r(|z|^2)$ are smooth and compactly
supported on $\R^2$; hence their radial Fourier transforms, and thus
$\mathcal H_cq$ and $\mathcal H_cr$, decrease faster than every power.
Set
\[
 a=\Lambda^{-1/2},\qquad b=\Lambda^{1/2},\qquad
 q_\Lambda(u)=q(u/a),\qquad r_\Lambda(u)=b^{-1}r(u/b).
\]
For large $\Lambda$ these functions are supported inside $(0,\Lambda)$,
with their prescribed coordinates exactly equal to $(1,0)$ and $(0,1)$.
Their raw norms are $O(\Lambda^{-1/4})$. The exact scaling identities
\[
 E_1[f(\cdot/a)]=a^{1/2}\tau_{\log a}E_1f,
 \qquad
 \mathcal H_c[f(\cdot/a)](v)=a(\mathcal H_cf)(av)
\]
show that their theta norms are also $O(\Lambda^{-1/4})$ and that,
for every $N$,
\begin{equation}\label{eq:hecke-lift-tail}
 \|(1-P_\Lambda)\mathcal H_cq_\Lambda\|_2+
 \|(1-P_\Lambda)\mathcal H_cr_\Lambda\|_2
                                  =O_N(\Lambda^{-N}).
\end{equation}
The corresponding absolute tail integrals have the same property.
Indeed the scaled lower integration limits are respectively
$a\Lambda=\Lambda^{1/2}$ and $b\Lambda=\Lambda^{3/2}$.

Put $A=P_\Lambda\mathcal H_cP_\Lambda\mathcal H_cP_\Lambda$
and $F=P_\Lambda\mathcal H_cP_\Lambda$, so $A=F^2$. For either
seed $f=P_\Lambda f$, the spectral theorem gives
\[
 \|(1-\Pi_\Lambda)f\|_2^2
 \leq\Lambda^{56}\langle(1-A)f,f\rangle
 =\Lambda^{56}\|(1-P_\Lambda)\mathcal H_cf\|_2^2.
\]
The fixed loss $\Lambda^{28}$ in norm is absorbed by
\eqref{eq:hecke-lift-tail}. Thus
$\|\Pi_\Lambda f-f\|_2=O_N(\Lambda^{-N})$ for every $N$.
The mean error is bounded by $\sqrt\Lambda$ times this norm.

Point evaluation needs a separate estimate. Hankel inversion for the
smooth seed gives
\[
 f(0)-(Af)(0)=c\int_\Lambda^\infty(\mathcal H_cf)(v)dv
                                           =O_N(\Lambda^{-N}).
\]
Since $(Fg)(0)=c\int_0^\Lambda g(u)du$ and $F$ is invertible on
the selected subspace,
\[
 |(Af)(0)-(\Pi_\Lambda f)(0)|
 \leq c\sqrt\Lambda\,
                 \|Ff-F^{-1}\Pi_\Lambda f\|_2.
\]
Commutation of $F$, $A$, and $\Pi_\Lambda$ bounds the last norm by
\[
 \|(1-\Pi_\Lambda)f\|_2+
 (1-\Lambda^{-56})^{-1/2}\|(A-1)f\|_2.
\]
Also $(1-A)f=P_\Lambda\mathcal H_c(1-P_\Lambda)\mathcal H_cf$.
Consequently the value error, as well as the mean error, is
$O_N(\Lambda^{-N})$. This uses the explicit band kernel rather than
an invalid point-evaluation estimate in fractional logarithmic Sobolev
space.

For any fixed $3/8<s<1/2$, the projected seeds have at most polynomial
$H_D^s$ norm by Proposition~\ref{prop:hankel-raw-estimates}; the
unprojected seeds have $H_D^s$ norms $O(\Lambda^{-1/4})$ by scaling.
Interpolation with the arbitrary-power raw error gives
\[
 \|\Pi_\Lambda f-f\|_{H_D^{3/8}}=O_N(\Lambda^{-N}).
\]
The graph bound \eqref{eq:hecke-graph} gives the same theta error.
The two projected coordinate columns form a matrix
$M_\Lambda=I_2+O_N(\Lambda^{-N})$. Multiplication by its inverse
produces $w_{1,\Lambda},w_{2,\Lambda}\in\operatorname{Ran}\Pi_\Lambda$
with coordinates exactly $(1,0)$ and $(0,1)$ and
\[
                  \|E_1w_{j,\Lambda}\|_2=O(\Lambda^{-1/4}).
\]
The constraint map therefore has rank two.

Since $J(\operatorname{Ran}\Pi_\Lambda\cap\ker I\cap\ker\operatorname{ev}_0)
=\operatorname{Ran}Q_\Lambda\oplus0$, the difference
$\mathcal Q_\Lambda-(Q_\Lambda\oplus0)$ is an orthogonal projection
of rank two. Its range is the graph over $\C^2$ of the operator
$X_\Lambda$ whose columns are $(1-Q_\Lambda)E_1w_{j,\Lambda}$.
Hence $\|X_\Lambda\|=O(\Lambda^{-1/4})$. The graph-projection
formula
\[
 P_{\operatorname{graph}X}=
 \begin{pmatrix}X\\I_2\end{pmatrix}
 (I_2+X^*X)^{-1}
 \begin{pmatrix}X^*&I_2\end{pmatrix}
\]
gives \eqref{eq:hecke-vertical-graph}; the difference from the vertical
projection has rank at most four and operator norm $O(\|X_\Lambda\|)$.
Finally raw unitary dilation by $e^t$ acts on the two coordinates by
$e^{-t/2}$ and $e^{t/2}$. Its theta action is $V(e^t)$, so its
integrated action is exactly the stated $R(h)$. Trace-norm convergence
against this bounded operator proves the last assertion.
\end{proof}

\subsection{Assembly for the auxiliary deep selection}
For each fixed character, Theorem~\ref{thm:hi-local-limit} identifies
the limit on compact frequency tests, and
Theorem~\ref{thm:hi-uniform-mass} extends it to every Schwartz test.
For $h\in C_c^\infty(\R)$ the function $H(i\tau)$ is Schwartz.
The exact compressed trace identity therefore gives
\[
 2Lh(0)-\Tr(S_LQ_{\chi,\Lambda}S_LV(h))
                   \longrightarrow\mathfrak B_\chi(h).
\]
The rank bound and Theorem~\ref{thm:hecke-deep-leakage} imply
\[
 \|Q_{\chi,\Lambda}-S_LQ_{\chi,\Lambda}S_L\|_1
 \leq2\sqrt{\operatorname{rank}Q_{\chi,\Lambda}}
              \|(1-S_L)Z_{\chi,\Lambda}\|_2
                         =O(\Lambda^{-13/8}).
\]
Thus the same limit holds for the uncompressed image projection.
Proposition~\ref{prop:hecke-principal-graph} proves the principal
polar correction. This completes the auxiliary deep-cutoff trace
formula. Theorem~\ref{nh:theorem} uses its upper bound and uniform
mass estimate to prove the fixed-threshold statement.

For the finite-order test algebra of Section~\ref{sec:complex-place},
only finitely many character blocks occur. Summing the sector identities
gives
\[
 \lim_{\Lambda\to\infty}
 \left(2\log\Lambda\,h(1)-\Tr(\mathcal Q_\Lambda R(h))
       +\int_\R h_1(t)(e^{-t/2}+e^{t/2})dt\right)
                 =\sum_{\chi\in F}\mathfrak B_\chi(h_\chi).
\]
Here $\mathcal Q_\Lambda$ is the direct sum of the sector projections;
its product with a finite-order test has finite rank. The limit is a
positive distribution on convolution squares. By the Hecke form of
Weil's criterion, it equals Weil's distribution on every test in these
sectors if and only if every finite-order Hecke function of $K$ has
its nontrivial zeros on the critical line. This implication uses neither
simple zeros nor a faithful representation of their full nilpotent jets.

\subsection{Comparison with a fixed concentration threshold}
\begin{proposition}[Comparison with the threshold $1/2$]
\label{prop:hecke-natural-upper}
Let $Q_\Lambda^{\mathrm{nat}}$ be the full-Gram theta-image projection
for $1_{[1/2,1]}(A_\Lambda)$, with the same principal constraints,
and let $Q_\Lambda^{\mathrm{deep}}$ denote the selection used above.
Let $\mu_\Lambda^{\mathrm{nat}}$ and $\mu_\Lambda^{\mathrm{deep}}$
be their compressed defect measures. Then
\[
 0\leq\mu_\Lambda^{\mathrm{nat}}\leq
                     \mu_\Lambda^{\mathrm{deep}}.
\]
In particular the natural measures have a uniform polynomial local-mass
bound, and for every nonnegative Schwartz function $\varphi$,
\[
 \limsup_{\Lambda\to\infty}\int\varphi\,d\mu_\Lambda^{\mathrm{nat}}
                              \leq\int\varphi\,d\mathfrak b_\chi.
\]
Writing $R_\Lambda=Q_\Lambda^{\mathrm{nat}}-Q_\Lambda^{\mathrm{deep}}$,
the natural compressed measures converge to $\mathfrak b_\chi$ on
Schwartz tests if and only if
\begin{equation}\label{eq:hecke-natural-extra-image}
 \|R_\Lambda S_LC_a\|_2\longrightarrow0
                           \qquad(a\in C_c^\infty(\R)).
\end{equation}
\end{proposition}
\begin{proof}
The deep raw subspace is contained in the natural one, also after
intersection with the two principal kernels. Their theta-image
subspaces are therefore nested, so $R_\Lambda$ is an orthogonal
projection and
\[
 0\leq S_L-S_LQ_\Lambda^{\mathrm{nat}}S_L
       \leq S_L-S_LQ_\Lambda^{\mathrm{deep}}S_L.
\]
Taking the Fourier diagonal gives the measure inequalities. The
uniform mass bound and the stated upper limit follow from the deep
theorem. Furthermore the exact trace identity gives
\[
 \int |a|^2\,d(\mu_\Lambda^{\mathrm{deep}}-
                          \mu_\Lambda^{\mathrm{nat}})
                 =\|R_\Lambda S_LC_a\|_2^2.
\]
Thus equality of the limits implies \eqref{eq:hecke-natural-extra-image}.
Conversely, that condition gives convergence on squared smooth compact
tests. Polarization gives all smooth compact tests, and the common
polynomial local-mass bound then gives every Schwartz test.
\end{proof}

This identifies the contribution of the additional theta-image
directions. Passing from the compression to the uncompressed trace
also requires
\[
 \Tr\bigl((Q_\Lambda^{\mathrm{nat}}-
                 S_LQ_\Lambda^{\mathrm{nat}}S_L)V(h)\bigr)\longrightarrow0.
\]
Neither assertion follows merely from the equality of the leading raw
ranks. The concentration calculation before arithmetic Gram normalization
does not by itself control these two image-space quantities.
Section~\ref{nh:section} proves the needed estimates by arithmetic
duality and mean-square summation.

